\documentclass[tikz,border=8pt]{standalone}
\usepackage{ctex,amsmath,amssymb}
\usetikzlibrary{arrows.meta,positioning,calc,shapes.geometric}
% Shared visual language for LFS architecture diagrams.
\RequirePackage{../../mymath}

\definecolor{LFSBlue}{HTML}{2563EB}
\definecolor{LFSBlueSoft}{HTML}{DBEAFE}
\definecolor{LFSGreen}{HTML}{15803D}
\definecolor{LFSGreenSoft}{HTML}{DCFCE7}
\definecolor{LFSOrange}{HTML}{C2410C}
\definecolor{LFSOrangeSoft}{HTML}{FFEDD5}
\definecolor{LFSRed}{HTML}{B91C1C}
\definecolor{LFSRedSoft}{HTML}{FEE2E2}
\definecolor{LFSPurple}{HTML}{7E22CE}
\definecolor{LFSPurpleSoft}{HTML}{F3E8FF}
\definecolor{LFSGray}{HTML}{475569}
\definecolor{LFSGraySoft}{HTML}{F1F5F9}

\tikzset{
  lfs picture/.style={
    font=\small,
    node distance=8mm and 12mm,
    >=Latex,
    every path/.style={line cap=round, line join=round}
  },
  block/.style={
    draw=LFSBlue,
    fill=LFSBlueSoft,
    rounded corners=2pt,
    line width=0.8pt,
    align=center,
    inner xsep=7pt,
    inner ysep=5pt,
    minimum height=10mm
  },
  compute/.style={
    block,
    draw=LFSPurple,
    fill=LFSPurpleSoft
  },
  storage/.style={
    block,
    draw=LFSGreen,
    fill=LFSGreenSoft
  },
  interface/.style={
    block,
    draw=LFSOrange,
    fill=LFSOrangeSoft
  },
  risk/.style={
    block,
    draw=LFSRed,
    fill=LFSRedSoft
  },
  neutral/.style={
    block,
    draw=LFSGray,
    fill=LFSGraySoft
  },
  decision/.style={
    diamond,
    draw=LFSOrange,
    fill=LFSOrangeSoft,
    line width=0.8pt,
    align=center,
    inner sep=2pt,
    aspect=2
  },
  group/.style={
    draw=LFSGray,
    rounded corners=3pt,
    dashed,
    line width=0.7pt,
    inner sep=6pt
  },
  arrow/.style={
    -{Latex[length=2.4mm,width=1.6mm]},
    draw=LFSGray,
    line width=0.9pt
  },
  data arrow/.style={
    arrow,
    draw=LFSBlue
  },
  control arrow/.style={
    arrow,
    draw=LFSOrange,
    dashed
  },
  residual/.style={
    arrow,
    draw=LFSGreen,
    rounded corners=4pt
  },
  label text/.style={
    font=\scriptsize,
    text=LFSGray,
    fill=white,
    inner sep=1.5pt
  },
  title/.style={
    font=\bfseries,
    text=LFSGray,
    align=center
  }
}


\begin{document}
\begin{tikzpicture}[lfs picture, node distance=6mm and 8mm]

% 上方：逻辑多维张量视图 (例如形状 [2, 3])
\node[title] at (5.5, 3.2) {逻辑二维视图：形状 $[2, 3]$，步长 $\text{Strides}=(3, 1)$};

\begin{scope}[shift={(3.2, 1.4)}]
  % 行 0
  \node[storage, text width=12mm, minimum height=8mm] (c00) {$(0,0)$\\$x_0$};
  \node[storage, text width=12mm, minimum height=8mm, right=2mm of c00] (c01) {$(0,1)$\\$x_1$};
  \node[storage, text width=12mm, minimum height=8mm, right=2mm of c01] (c02) {$(0,2)$\\$x_2$};

  % 行 1
  \node[storage, text width=12mm, minimum height=8mm, below=2mm of c00] (c10) {$(1,0)$\\$x_3$};
  \node[storage, text width=12mm, minimum height=8mm, right=2mm of c10] (c11) {$(1,1)$\\$x_4$};
  \node[storage, text width=12mm, minimum height=8mm, right=2mm of c11] (c12) {$(1,2)$\\$x_5$};

  \node[label text, left=3mm of c00] {行 0};
  \node[label text, left=3mm of c10] {行 1};
\end{scope}

% 下方：物理显存线性排布 (Row-Major 连续存储)
\node[title] at (5.5, -0.2) {物理显存一维连续排布（地址偏移 $\text{Offset} = i \times 3 + j \times 1$）};

\begin{scope}[shift={(0.8, -1.3)}]
  \foreach \i/\val/\idx in {0/$x_0$/0, 1/$x_1$/1, 2/$x_2$/2, 3/$x_3$/3, 4/$x_4$/4, 5/$x_5$/5} {
    \node[block, text width=11mm, minimum height=9mm] (m\i) at (\i*1.6, 0) {\val\\\scriptsize 偏移 \idx};
  }
\end{scope}

% 映射虚线箭头
\draw[data arrow, dashed] (c00.south) -- (m0.north);
\draw[data arrow, dashed] (c02.south) -- (m2.north);
\draw[data arrow, dashed] (c10.south) -- (m3.north);
\draw[data arrow, dashed] (c12.south) -- (m5.north);

% 底部注释
\node[label text] at (5.5, -2.4) {转置变换 $\mat{X}^{\mathrm{T}}$ 仅修改步长为 $(1, 3)$（变为非连续张量）；调用 \texttt{.contiguous()} 触发显存重新分配与物理深拷贝};

\end{tikzpicture}
\end{document}
